\section{Attaching the funded service to the selected history}
\label{sec:original-history-portfolio-attachment}

The finite-volume service is useful only if its data belong to the history
selected for the public theorem. We verify that binding here. The public
construction selects one run and applies the funded continuation to that
run's own history, schedule, transitions, and finite-volume data. No cutoff,
period, or extraction order is used to select a second history. The common
schedule and its numerical prices are organized in
Section~\ref{sec:common-funded-history}; here we check the remaining
same-witness interface.

Let $R$ denote the selected run. Write $Y_R$ for its stored finite
renormalization-group history and $\mathcal P_R$ for the separately stored
scale campaign.
The distinction matters: the history used below is
$\operatorname{trail}(Y_R)$, whereas the campaign supplies the scale data.
Let $t_R(d)$ be the reference level at grade $d$, let $S_R(d,N,p)$ be the
transfer datum, and let $D_R(N,p,m)$ be the selected transition record. We
write
\[
 C_R=\bigl(\operatorname{trail}(Y_R),\ t_R,\ S_R\bigr)
\]
for the family with its fixed holonomy, label frames, and local sections.
This notation names the family already stored in the run; it does not define
a new history.

\begin{proposition}[Same-history attachment]
\label{prop:ex-original-history-attachment}
Suppose the selected run $R$ carries its funded continuation certificate.
Applying that certificate to the data stored in $R$ gives a package on the
same family $C_R$ with the following properties:
\begin{enumerate}
\item the reference-level condition
\[
 d+3\leq J(t_R(d))\qquad(d\geq0);
\]
\item one nesting family, one scale schedule, one transition family, and
   one transfer family, such that the selected transition at $(N,p,m)$ is
   available for $m<N$, and, for every $d,N,p$ with $t_R(d)\leq N$, the
   finite-volume transfer identity holds using those same data;
\item the funded nonlinear-response and adjacent-scale estimates, together
   with the selected service on exactly $C_R$;
\item the centered service, including its height estimates and terminal
   forced-profile relation, as well as the finite-volume bounds and marked
   portfolio conclusions recorded for $C_R$.
\end{enumerate}
The first three items and the finite-volume identities hold simultaneously
for all cutoffs and periods in their stated domains. For each fixed grade
and section in the purchased portfolio, the terminal marked estimate is
conditional on the reserved addresses and fee, the local forced-family
attachment, the paid selected-history record, and the served-representation
input used in Propositions~\ref{prop:ex-marked-invitation} and
\ref{prop:actual-marked-figment}. Under these conditions on the same
$C_R$, it holds for all sufficiently large cutoffs, uniformly in the period;
its threshold may depend on that grade and section. It is a separate local
estimate, not a field included implicitly in the centered service.
\end{proposition}

\begin{proof}
We first identify the history-bearing data. The selected run stores $Y_R$,
the schedule $t_R$, a transition family, the gradewise regularity data, and
the finite-volume transfer witnesses. Its underlying family is
\[
 C_R=\bigl(\operatorname{trail}(Y_R),t_R,S_R\bigr),
\]
with the fixed holonomy, label frames, and local-section map. The funded
continuation is a function of precisely these stored objects: the
configuration and its nesting data, the schedule and its opening
inequalities, the transfer family and transition records, and the
regularity and finite-volume transfer witnesses. Its output therefore has
the family formed from those same arguments. That family is $C_R$ itself;
there is no intervening identification map and no reselected label system.

Here is why the quantifiers in that application are uniform. The public
selection fixes the history $Y_R$ and a nesting family for it.
Independently, it chooses the reference schedule on that selected run so
that $d+3\leq J(t_R(d))$ for each grade. On this fixed history it constructs
the transition branches and their records once, and then defines one
transfer family on all triples $(d,N,p)$. That family has the required
regularity properties; for every $d,N,p$ with $t_R(d)\leq N$, its
finite-volume witness is built from the already chosen nesting family and
transition records. The witness at $(d,N,p)$ may depend on those indices,
but the history, nesting family, schedule, transitions, and transfer family
do not. Thus the simultaneous quantifiers are part of the selected run and
are not obtained by diagonal extraction or by making new global choices
after fixing $(N,p)$.

For the transfer witnesses, if $t_R(d)>N$ the construction uses a default
map. If $t_R(d)\leq N$, it composes along the selected transitions and
obtains the finite-volume identity and accounting estimate. This explains
the cutoff in item 2 and shows that both cases use the original $Y_R$.

The funded continuation supplies two quantitative inputs to the centered
construction: bounds on the delayed nonlinear response and its tariff, and
uniform estimates for each allowed adjacent-scale step. The latter are the
scale inequalities proved in
Section~\ref{ex-centered-geyser-scale-estimates}. That section applies them
to the selected adjacent levels and block size, and supplies a common
monotone threshold. It does not change the selected history, nor does it
by itself establish the nonlinear-response bounds.

The centered service is assembled from these data as follows. The escape
estimate uses the same nesting family, nonlinear-response bounds, and all
selected transitions. Its room and bank estimates use the budget fields of
the selected factor. The adjacent-scale inequalities give the required
scale comparison at each selected step. For heights at grade $d$ and
reference level $k$, the local residence estimate at $k$ gives the later
scale bounds, the one-cover footprint estimate, and a uniform camera
bound. The five numerical implications used in this conversion are proved
in Proposition~\ref{prop:ex-geyser-service-fields}; the eventual estimates
and finite-prefix maximum are in
Theorem~\ref{thm:ex-selected-geyser-steam}. The maximum is over the finitely
many grades below each target stage and its two adjacent stages, so it gives
a common threshold for all later uses at that stage.

At the terminal endpoint, the finite-volume transfer witness gives the top
transfer identity and its forced-profile relation. The centered service's
terminal estimate is this same relation, using the domination and rise
identities attached to the selected run. At an intermediate scale, the
transfer witnesses and selected transitions feed the local transition
estimate: specializing the tariff inequality to the zero-grade section
gives the required relay bounds. The separate growth and jump estimates
follow from positivity and the selected scale payment. The three-card
coefficient, exponential-entry, and forced-profile identities are proved
in Section~\ref{sec:three-port-history}; we use those identities here
rather than constructing a second interpolation or repeating their price
proof.

The funded continuation returns the centered service together with its
finite-volume bounds and marked portfolio. Its local transition estimates
use the same nonlinear-response bounds, adjacent-scale estimates, nesting
family, transition records, finite-volume witnesses, and schedule just
described. Both the core centered-service assembly and its stronger funded version
use the supplied configuration, schedule, transfer data, and transition
family. Neither creates a second history.

Finally, the terminal marked-field estimate is attached only after $C_R$
has been fixed. The reserved addresses of the terminal representation
are not constructed there: they are a property of the selected history,
recorded for every cutoff and period in the nesting receipt stored in the
run. This receipt is an input of the present section; it originates in the
selection of the history from the exact startup family, and no result of
this group of sections proves it. Given the reserved addresses and the
lawful terminal representation,
Proposition~\ref{prop:terminal-marked-voice} constructs the terminal
marked coefficients on this same family, with their density and
forced-profile identities, provided its fee inequality holds; that
inequality is the terminal fee purchased in
Section~\ref{sec:common-observable-purchases}. The quantitative local
bounds are proved in Section~\ref{sec:ex-marked-field-local}, and their
eventual uniformity for each fixed grade and section in
Proposition~\ref{prop:ex-marked-solaris}. The local estimate requires soundness,
reference-level eligibility, fit of the union of the stencils, nonnegative
credit, a keeper, the regional card-times-allowance bound, cover at most
$81$, the square-root allowance-times-storm bound, and the stated positive
terminal weight and nonnegative reserve. The local estimates supply the
numerical bounds for all sufficiently large cutoffs for each fixed grade
and section in the purchased portfolio, uniformly in the period, subject to
the stated schedule and matching inputs. In addition, the local
forced-family attachment must be supplied for each relevant correction,
on the terminal representation of this same $C_R$, uniformly in the
period. The paid selected-history record and served-representation input
are also required for the root-test conclusion; their joint production on
$C_R$ is not proved here. Subject to these inputs, the marked density and
forced-profile identity give the terminal estimate on $C_R$.
Throughout, its history, schedule, transfer family, and terminal
representation remain those of $C_R$. No uniqueness assertion for a
terminal field is needed.
\end{proof}

\begin{remark}[Scope of the attachment]
The proposition proves the same-history attachment once the selected run's
funded continuation certificate is available. The common funded-history
construction is responsible for producing that certificate and its
quantitative prices; this section does not replace those purchases by the
name of a record. It also does not use the separate grantor route as
evidence for the public caller. The public construction applies its funded
continuation to the data stored in the selected run.

Two points of bookkeeping are needed by later sections. First, the
certificate is a universally quantified statement: it applies to every
history of the same alloy with a nesting receipt, every schedule obeying
the reference-level condition, and all transfer and transition families
satisfying the making, climbing and transfer conditions. The run stores
this certificate, and the public run stores in addition a stronger form
with further conclusions; both are specialized to the same nine stored
objects, namely the history, its nesting receipt, the schedule and its
reference-level condition, the transfer family, the transition family,
and the making, climbing and transfer receipts. Second, the certificate
asserts the existence of the marked portfolio, that is, of one local
section and one marked family for each grade together with the
gradewise coefficient maps and the presentation entering the terminal
marked estimate; it does not name them. A later argument must fix one
such witness once, before any cutoff, period or extraction order is
chosen, and use it throughout. Reopening the existence statement at a
later point may produce a different portfolio and is not a same-object
argument.
\end{remark}

\subsection{Budget and domain of the stored service}
\label{sec:ex-stored-service-budget}

Write $W(j)$ for the worth fixed in the purchase, $c_j$ for the
column count of stage $j$, and $\ell_j$ for its chosen block side.
At level $i$ put $C_i=c_{J(i)}$ and $L_i=\ell_{J(i+1)}$.
Thus $C_i$ is the column count denoted by $\Gamma_i$ in
Theorem~\ref{thm:ex-selected-run-service}; it is unrelated to the
running-coupling increment. Notice that the level-$i$ block is
chosen at the \emph{entrance} stage $J(i+1)$.

\begin{proposition}[Room and bank on the purchased levels]
\label{prop:ex-stored-service-budget}
The stage purchase can be made so that every allowed schedule has,
at every level $i\geq0$,
\[
 64 W(J(i)+1)^2\le C_i,\qquad
 100000 W(J(i+1)+1)^6(1+C_{i+1})^2\le L_i^{1/8}.
\]
In particular the same inequalities hold for the selected run,
independently of grade, cutoff and period.
\end{proposition}

\begin{proof}
After fixing worth and the orders, choose the stage column count
as a running maximum of the room request $64W(j+1)^2$, the
factorial of the stage derivative order and the other pre-length
column requests. An integer count can be used by taking ceilings.
In particular the profile already pays the room request before any
length is chosen. With these columns fixed, let
\[
 A_j=100000W(j+1)^6(1+c_j)^2,\qquad
 n_j=\max\{1,\lceil A_j^8\rceil\}.
\]
Include $n_j$ in the stage length floor before selecting $\ell_j$.
This is compatible with all other lower floors: take their maximum
and, if nondecreasing floors are required, take the running maximum
of that maximum. Any subsequent increase preserves the two
purchases. Since $\ell_j\ge n_j$ and the eighth root is increasing
on the nonnegative reals, $A_j\le\ell_j^{1/8}$. Substitute $j=J(i)$
in the column inequality and $j=J(i+1)$ in the length inequality.
The same substitutions apply at $i=0$. The later entrance delays
and the selection of finite-volume histories do not change the
stage columns or these lengths.
\end{proof}

For the common C02 and C12 interface use the direct specialization
of the funded continuation in
Proposition~\ref{prop:ex-original-history-attachment}, with family
$C_R$, reference map $t_R$, transfer data $S_R$ and doors $D_R$.
The stronger funded service is on precisely these objects. Its
centered part restricts to each reference level $k$ by retaining
the transition receipts with $k\le i<N$ and the terminal receipt
at $N$. The stock clause for grade $d$ is needed only for
$t_R(d)\le k\le N$, and the music clause for $t_R(d)\le N$ with
doors $t_R(d)\le i<N$. The restriction makes sense even when
$k=N$, with no intervening door. The centered birth estimate uses
the stricter domain $k<N$. All period and cutoff indices range
over entries of this one run. The local analytic attachments
needed for retention, music, the paid local coupon and the uniform
birth estimate are separate from this projection and the numerical
budget proposition.

There is also a campaign construction with its own chosen
finite-volume history, transfers and doors. In the stamped original
campaign the history is $Y_R$ and the transfers are $S_R$, so
forming its family gives $C_R$ by substitution. The campaign
service uses the recorded gate of its campaign step for $i\ge1$
and its $k$-restricted retention is asserted for $k\ge1$.
At $i=0$ its door is a default. The common interface above uses
$D_R$ directly, including at level zero, and needs no identification
of campaign gates with these doors.

For this direct interface put $x_i=\widehat\beta_i$ and
$r_i=N_{J(i+1)}$, where $N_j$ is the stage occupancy bound.
The marched weight multiplier is $\phi_i=x_{i+1}^{-1}$ and the
probe price is $\delta_i=x_{i+1}^{-2}$; both use the fine coupling.
The grade height is
$\zeta_{d,i}=x_i^{-(\deg F_d+4)}$ and the response waybill
$\psi_i$ is twice the title cost at level $i+1$ multiplied by
the exponential of that level's support exponent.
The terminal-bank reserve is $u_i=x_i^{-1/2}/4$.
In particular it is neither the marched multiplier nor the
adjacent coupon rate. Proposition~\ref{prop:selected-centered-retention}
gives the separate terminal endowment from the paid delay at $J(N+1)$:
for $C_N^{\rm end}=w_N^{\rm f}+r_N\mathsf s(B_N)u_N$,
the bounds are $(2B_N+1)^4 C_N^{\rm end}\le3$ and
$e^{C_N^{\rm end}}-1\le1$, where $w_N^{\rm f}$ is the forced
allowance. These numerical bounds require no door beyond the cutoff.
